\documentclass[12pt]{article}
\usepackage[utf8]{inputenc}
\usepackage[T1]{fontenc}
\usepackage{hyperref}
\title{Filomena Mata-Coahuitl\'an Totonac Language Resources for K-12 Curriculum Development}
\author{Compiled by Claude Sonnet 5 (claude-sonnet-5) for Andr\'as Kornai}
\date{}
\begin{document}
\maketitle

\section{Introduction}
Filomena Mata-Coahuitl\'an Totonac (also written Filomeno Mata Totonac; ISO 639-3 tlp) is a Totonacan language of Veracruz, Mexico. Estimates range from about 15,000 speakers (Wikipedia, citing 2000) to about 17,000 (Global Recordings Network). A 2024 study in \emph{Acta Antropol\'ogica} on Filomeno Mata families reports strong transmission through grandparents but no language policy supporting bilingualism at home. Its place within the Northern or Sierra branch of Totonacan is debated.

\section{Primary Educational Resources}
The searches found no school materials specific to this variety.

\section{Linguistic and Reference Materials}
\begin{itemize}
\item Wikipedia entry: \url{https://en.wikipedia.org/wiki/Filomeno_Mata_Totonac}
\item Wikipedia entry on the Totonac languages: \url{https://en.wikipedia.org/wiki/Totonac_languages}
\item Wikitongues language page: \url{https://wikitongues.org/languages/tlp/}
\item Numerals data, Max Planck Institute for Evolutionary Anthropology: \url{https://lingweb.eva.mpg.de/channumerals/Totonac-Filomena-Mata-Coahuitlan.htm}
\item Coahuitl\'an Totonac Questionnaire in the AILLA archive at UT Austin: \url{https://islandora-ailla.lib.utexas.edu/islandora/object/ailla%3A137003}
\item 2024 \emph{Acta Antropol\'ogica} article on Filomeno Mata families: \url{https://www.scielo.org.mx/pdf/aa/v59n1/2448-6221-aa-59-01-139.pdf}
\end{itemize}
Teresa McFarland's 2009 Berkeley dissertation on the phonology and morphology of the language was mentioned in the search summary. No URL for it appeared, so none is given.

\section{Religious and Community Resources}
No scripture or community organisation was found in the searches.

\section{Digital and Audio Resources}
\begin{itemize}
\item Global Recordings Network language page (I did not check what recordings it lists): \url{https://globalrecordings.net/en/language/tlp}
\end{itemize}

\section{Curriculum Development Context}
The language is vigorous and is passed on at home, which favours a native-language curriculum. No existing school materials were found, so these sources are mainly linguistic documentation.
\end{document}